\documentclass[
	fontsize=12pt,
	paper=a4,
	numbers=noenddot,
	listof=totoc,
	listof=entryprefix,
	listof=nochaptergap,
	bibliography=totoc,
	parskip=half,
	openany
]{scrbook}

\usepackage{cmap}


\usepackage[utf8]{inputenc}
\usepackage[T1]{fontenc}
\usepackage{lmodern}
\usepackage{ngerman}


\usepackage{scrhack}
\usepackage[table]{xcolor}
\usepackage{chngcntr}
\usepackage{enumitem}


\usepackage{amsthm}


\usepackage{lipsum}
\usepackage{unicode-math}


\usepackage{mathtools}


\usepackage[
	colorlinks=false,
	linkcolor=black,
	citecolor=black,
	filecolor=black,
	urlcolor=black,
	bookmarks=true,
	bookmarksopen=true,
	bookmarksopenlevel=3,
	bookmarksnumbered,
	plainpages=false,
	pdfpagelabels=true,
	hyperfootnotes
]{hyperref}


\renewcommand{\familydefault}{\sfdefault}
\usepackage[onehalfspacing]{setspace}
\raggedbottom

\usepackage{anyfontsize}

\addtokomafont{chapter}{\sffamily\large\bfseries}
\addtokomafont{section}{\sffamily\normalsize\bfseries}
\addtokomafont{subsection}{\sffamily\normalsize\mdseries}
\addtokomafont{caption}{\sffamily\normalsize\mdseries}

\usepackage{relsize}
\renewcommand*{\UrlFont}{\ttfamily\smaller\relax}


\usepackage[
	bindingoffset=0cm,
	inner=2.5cm,
	outer=2.5cm,
	top=3cm,
	bottom=2cm
]{geometry}

\usepackage[
	headsepline,
	plainheadsepline
]{scrlayer-scrpage}
\clearpairofpagestyles
\addtokomafont{pagehead}{\normalfont}
\ohead*{\thepage}
\ihead*{\leftmark}


\usepackage{graphicx}
\graphicspath{{pictures/}}
\usepackage{float}

\renewcommand{\thefigure}{\arabic{figure}}
\counterwithout{figure}{chapter}

\usepackage{wrapfig}


\usepackage{booktabs}
\usepackage{multirow}
\usepackage{tabularx}

\renewcommand{\thetable}{\arabic{table}}
\counterwithout{table}{chapter}


\usepackage{listings}
\usepackage{beramono}

\renewcommand{\lstlistlistingname}{List of Code Listings}
\renewcommand{\lstlistingname}{Code Listing}
\newcommand{\listoflolentryname}{\lstlistingname}

\definecolor{codegreen}{rgb}{0,0.6,0}
\definecolor{codegray}{rgb}{0.5,0.5,0.5}
\definecolor{codepurple}{rgb}{0.5,0,0.33}
\definecolor{codepurblue}{rgb}{0.16,0.0,1.0}
\definecolor{backcolour}{rgb}{0.95,0.95,0.92}

\lstdefinestyle{codestyle}{
	backgroundcolor=\color{backcolour},
	commentstyle=\color{codegreen},
	keywordstyle=\bfseries\color{codepurple},
	numberstyle=\tiny\color{codegray},
	stringstyle=\color{codepurblue},
	basicstyle=\scriptsize\ttfamily,
	breakatwhitespace=false,
	breaklines=true,
	captionpos=b,
	keepspaces=true,
	numbers=left,
	numbersep=5pt,
	showspaces=false,
	showstringspaces=false,
	showtabs=false,
	tabsize=2,
	escapeinside={(*@}{@*)}
}
\lstset{style=codestyle, numberbychapter=false}



\KOMAoptions{toc=chapterentrydotfill}
\addtokomafont{chapterentry}{\normalfont}
\setuptoc{toc}{totoc}

\DeclareTOCStyleEntry[beforeskip=0cm]{chapter}{chapter}
\DeclareTOCStyleEntry[beforeskip=0cm]{section}{section}
\DeclareTOCStyleEntry[beforeskip=0cm]{default}{subsection}

\BeforeStartingTOC[lof]{\def\autodot{:}}
\BeforeStartingTOC[lot]{\def\autodot{:}}
\BeforeStartingTOC[lol]{\def\autodot{:}}


\iffalse
\usepackage{csquotes}
\setlength\bibitemsep{.5\baselineskip}
\setcounter{biburlnumpenalty}{9000}
\setcounter{biburllcpenalty}{9000}
\setcounter{biburlucpenalty}{9000}

\fi

\usepackage[printonlyused]{acronym}


\usepackage[title,titletoc]{appendix}

\newcommand{\appendixchapter}[1]{
\cleardoublepage
\pagenumbering{arabic}
\renewcommand{\thepage}{\thechapter-\arabic{page}}
\chapter{#1}
}

\newcommand{\monthlyreport}[2]{
\section{#1}
\centering
\includegraphics[trim=55 35 55 35,clip,width=1\textwidth]{#2}
\clearpage
}



\newtheorem{beispiel}{Beispiel}

\newtheorem{these}{These}

\newtheorem{definition}{Definition}


\usepackage[nameinlink, noabbrev]{cleveref}


\lstset{
	extendedchars=true,
	inputencoding=utf8,
	literate=
		{ä}{{\"a}}1
		{ö}{{\"o}}1
		{ü}{{\"u}}1
		{ß}{{\ss}}1
		{Ä}{{\"A}}1
		{Ö}{{\"O}}1
		{Ü}{{\"U}}1
}

\begin{document}

\begin{titlepage}
	\pagestyle{empty}

	\begin{flushright}
		\begin{figure}[ht]
			\flushright
			\includegraphics[height=2cm]{../../for_latex/hfu}
		\end{figure}
	\end{flushright}

	\begin{center}
		\vspace{3cm}

		{\fontsize{22}{22} \selectfont \textbf{Blatt 8}}\\[5mm]
		{\fontsize{18}{18} \selectfont Web-Security und Intrusion Prevention}

		\vspace{12cm}

		{\fontsize{14}{14} \selectfont Luis Staudt}
	\end{center}
\end{titlepage}


\chapter*{Vorbemerkung}

Ein Teil dieser Aufgabe setzt externe Infrastruktur voraus, die lokal nicht
reproduzierbar ist: eine bwCloud-VM mit öffentlicher IP (Login über
HFU-Account), einen öffentlichen DNS-Namen über duckdns.org und ein öffentlich
validiertes Let's-Encrypt-Zertifikat, dessen ACME-Challenge von aussen
erreichbar sein muss. Dieser Teil wurde nicht ausgeführt.

Der zentrale Lerninhalt, der MITM-Angriff per ARP-Spoofing und der Schutz durch
HTTPS, wurde dagegen in einem lokalen, self-contained Docker-Lab durchgeführt
(Kapitel \glqq Lokale Reproduktion\grqq{}), mit einem lokalen Login-Server
anstelle der bwCloud-Seite. HSTS und die Let's-Encrypt-Validierung werden
zusätzlich erklärt. Das Lab liegt reproduzierbar unter \texttt{lokal-lab/}
(\texttt{run-mitm.sh}).

\chapter*{Aufgabe 1: Die Anwendungen}

Die lokale Angriffsumgebung besteht aus drei Docker-Containern:

\begin{itemize}
	\item \textbf{browser} -- das Opfer. Ein Firefox, der über VNC
	      beziehungsweise eine Web-Oberfläche (\url{http://10.0.0.10:5800/})
	      bedient wird und sich über den Router mit dem Internet verbindet.
	\item \textbf{router} -- stellt die Internet-Verbindung her
	      (NAT/MASQUERADE, \texttt{net.ipv4.ip\_forward=1}). Default-Gateway ist
	      \texttt{90.0.255.254}.
	\item \textbf{attacker} -- der Angreifer mit den Werkzeugen
	      dsniff/arpspoof, tcpdump, scapy und bettercap. Führt den MITM-Angriff
	      durch.
\end{itemize}

Der Webserver (Immich) läuft getrennt davon auf der bwCloud-VM.

\chapter*{Aufgabe 2: VM und Webanwendung (bwCloud)}

Die folgenden Schritte wurden auf der externen bwCloud-Infrastruktur ausgeführt,
nicht im lokalen Lab:

\begin{itemize}
	\item SSH-Schlüssel erzeugen (\texttt{ssh-keygen}, mit Passwort gesichert);
	      der private Schlüssel bleibt geheim, \texttt{id\_rsa.pub} wird in die
	      VM importiert. Der Passwort-Login ist deaktiviert, da Passwörter leicht
	      angreifbar sind.
	\item Debian-13-VM (m1.medium, public-v4) in bwCloud erstellen, per SSH
	      verbinden (\texttt{ssh debian@<IP>}), Updates einspielen, Docker
	      installieren und den Benutzer \texttt{debian} zur \texttt{docker}-Gruppe
	      hinzufügen.
	\item Domainnamen über duckdns.org registrieren und die statische VM-IP
	      eintragen.
	\item Immich per docker-compose installieren, Port vorübergehend auf 80
	      stellen und einen Testbenutzer anlegen, mit Pseudonym und
	      Wegwerf-Passwort, da die Verbindung zunächst unverschlüsselt und
	      mitlesbar ist.
\end{itemize}

\chapter*{Aufgabe 3: MITM-Angriff per ARP-Spoofing}

ARP kennt keine Authentifizierung. Der Angreifer sendet gefälschte ARP-Antworten
und bringt so Opfer und Gateway dazu, seine MAC-Adresse für die jeweils andere IP
einzutragen. Dadurch läuft der gesamte Verkehr über den Angreifer
(Man-in-the-Middle).

Im attacker-Container wird zunächst das Default-Gateway gesetzt, da andernfalls
die Internet-Verbindung des Opfers abbricht:

\begin{lstlisting}[language=bash]
docker exec -it attacker bash
route add default gw 90.0.255.254
\end{lstlisting}

Das ARP-Spoofing erfolgt in beide Richtungen: dem Opfer 90.0.0.1 gegenüber als
Gateway und dem Gateway gegenüber als Opfer:

\begin{lstlisting}[language=bash]
# gegenueber dem Opfer als Gateway auftreten
arpspoof -i eth0 -t 90.0.0.1     90.0.255.254
# gegenueber dem Gateway als Opfer auftreten
arpspoof -i eth0 -t 90.0.255.254 90.0.0.1
\end{lstlisting}

Da Immich zunächst über HTTP läuft, stehen die Login-Daten im Klartext im Paket.
tcpdump zeigt die Paketinhalte an:

\begin{lstlisting}[language=bash]
tcpdump -i eth0 -A -n 'tcp port 80'
\end{lstlisting}

Loggt sich das Opfer nun auf der Webseite ein, erscheint das im Klartext
übertragene Passwort im Mitschnitt (POST-Body des Login-Requests). Danach werden
\texttt{arpspoof} und \texttt{tcpdump} beendet.

\chapter*{Aufgabe 4: Schutz mit Verschlüsselung und Zertifikaten}

Immich wird auf Port 2283 zurückgestellt und ist nur noch intern erreichbar.
Davor kommt der \emph{NginX Proxy Manager} (NPM), der über Port 80/443 nach
aussen spricht und per Let's Encrypt ein Zertifikat für die duckdns-Domain
anfordert. Zusätzlich werden \emph{Force SSL} und \emph{HSTS} aktiviert. Port 81
(NPM-Admin) muss in der bwCloud-Firewall freigegeben werden.

\section*{Bedeutung von HSTS}

HSTS (HTTP Strict Transport Security) ist ein Response-Header
(\texttt{Strict-Transport-Security: max-age=...}), mit dem der Server dem Browser
mitteilt, die Seite für die angegebene Dauer ausschliesslich über HTTPS
aufzurufen. Daraus ergeben sich folgende Konsequenzen:

\begin{itemize}
	\item Der Browser ersetzt \texttt{http://}-Aufrufe der Domain intern durch
	      \texttt{https://}, noch bevor eine Anfrage das Netz verlässt.
	\item Ein erzwungenes Downgrade auf HTTP (SSL-Stripping) wird verhindert.
	\item Bei Zertifikatsfehlern ist kein \glqq Trotzdem fortfahren\grqq{} mehr
	      möglich.
\end{itemize}

HSTS schliesst damit das Zeitfenster der ersten, noch unverschlüsselten Anfrage,
mit \texttt{preload} sogar vollständig.

\section*{Wie Let's Encrypt den Besitz des Servers prüft}

Let's Encrypt nutzt das ACME-Protokoll mit einer
Domain-Validation-Challenge. Der Antragsteller muss nachweisen, dass er die
Domain kontrolliert:

\begin{itemize}
	\item \textbf{HTTP-01:} Let's Encrypt gibt ein Token vor; der Server muss es
	      unter \texttt{http://<domain>/.well-known/acme-challenge/<token>}
	      bereitstellen. Let's Encrypt ruft diese URL über den öffentlichen
	      DNS-Namen ab. Nur wer den Server hinter der Domain kontrolliert, kann
	      die Datei dort ablegen.
	\item \textbf{DNS-01:} Alternativ wird ein vorgegebener Wert als TXT-Record
	      (\texttt{\_acme-challenge.<domain>}) im DNS hinterlegt.
\end{itemize}

Ein Angreifer ohne Kontrolle über Domain beziehungsweise Server kann die
Challenge nicht erfüllen und erhält daher kein gültiges Zertifikat. NPM übernimmt
diesen Ablauf automatisch.

\section*{MITM nach der Umstellung}

Wird der Angriff bei aktivem HTTPS wiederholt, sieht der Angreifer nur noch
TLS-verschlüsselten Verkehr; das Passwort ist nicht mehr auslesbar. Zwar könnte
er die Pakete weiterhin umleiten, ohne den privaten Schlüssel des Servers aber
den Inhalt nicht entschlüsseln. Der Versuch, ein eigenes Zertifikat
unterzuschieben, scheitert an der Zertifikatsprüfung des Browsers und zusätzlich
an HSTS.

\chapter*{Lokale Reproduktion (Docker)}

Da die bwCloud-Seite extern ist, wurde der Angriff in einem isolierten
Docker-Netz (\texttt{192.168.80.0/24}) nachgestellt. Es kamen drei Container zum
Einsatz: \texttt{server} (\texttt{.10}, Login-Webseite), \texttt{victim}
(\texttt{.20}) und \texttt{attacker} (\texttt{.66}).

\section*{Aufgabe 3 lokal: ARP-Spoofing + Klartext-Passwort}

Der Angreifer aktiviert IP-Forwarding und fälscht per \texttt{arpspoof} die
ARP-Zuordnung in beide Richtungen:

\begin{lstlisting}[language=bash]
sysctl -w net.ipv4.ip_forward=1
arpspoof -i eth0 -t 192.168.80.20 192.168.80.10   # dem Opfer: ich bin der Server
arpspoof -i eth0 -t 192.168.80.10 192.168.80.20   # dem Server: ich bin das Opfer
\end{lstlisting}

Danach zeigt die ARP-Tabelle des Opfers für die Server-IP die MAC des Angreifers,
sodass der Verkehr nun über den Angreifer läuft:

\begin{lstlisting}[language=bash]
# victim /proc/net/arp
192.168.80.10   ...   56:24:2a:f5:f7:52   *   eth0   # = MAC des Angreifers
\end{lstlisting}

Loggt sich das Opfer nun per HTTP ein, liest der Angreifer das Passwort im
Klartext aus dem Mitschnitt (\texttt{images/mitm\_http.pcap}):

\begin{lstlisting}[language=bash]
tcpdump -nr mitm.pcap -A | grep password=
-> user=alice&password=SuperGeheim123
\end{lstlisting}

\section*{Aufgabe 4 lokal: mit HTTPS ist das Passwort geschützt}

Derselbe Angriff richtet sich gegen einen TLS-Server (Port 443, selbstsigniertes
Zertifikat). Das Opfer loggt sich jetzt über HTTPS ein. Der Angreifer sieht zwar
die Pakete, aber nur TLS-verschlüsselt:

\begin{lstlisting}[language=bash]
# Vorkommen des Passworts im jeweiligen Mitschnitt:
grep -c SuperGeheim  (HTTP-Mitschnitt mitm.pcap)      -> 4   (lesbar)
grep -c SuperGeheim  (HTTPS-Mitschnitt mitm_tls.pcap) -> 0   (nicht lesbar)
\end{lstlisting}

Ohne Verschlüsselung ist das Passwort trivial mitlesbar, mit HTTPS nicht. Dies
deckt sich mit dem erwarteten Verhalten auf der bwCloud. Die pcap-Dateien liegen
unter \texttt{images/} bei.

\section*{Zertifikatsausstellung per ACME (lokal mit Pebble)}

Ein echtes Let's-Encrypt-Zertifikat erfordert eine öffentliche Domain, deren
ACME-Challenge von aussen erreichbar ist, was lokal nicht möglich ist.
Stattdessen wurde \emph{Pebble}, der offizielle Test-ACME-Server von Let's
Encrypt, verwendet. Pebble spricht dasselbe ACME-Protokoll (Order, HTTP-01,
Ausstellung) wie die produktive CA; lediglich die ausstellende CA ist eine
Test-CA. So lässt sich der vollständige Ablauf real durchführen.

Die Ausstellung erfolgt über \texttt{acme.sh} im Standalone-Modus. Pebble
validiert die HTTP-01-Challenge, indem es
\url{http://<domain>:5002/.well-known/acme-challenge/<token>} abruft, also genau
die Datei, die der ACME-Client während der Ausstellung bereitstellt:

\begin{lstlisting}[language=bash]
acme.sh --issue --standalone --httpport 5002 -d immich.websec.lab \
        --server https://pebble:14000/dir
-> Order -> HTTP-01 challenge -> finalize -> Cert success.
\end{lstlisting}

Das ausgestellte Zertifikat zeigt im Auszug (\texttt{images/acme\_cert.txt}):

\begin{lstlisting}[language=bash]
issuer=CN=Pebble Intermediate CA 1a139c
X509v3 Subject Alternative Name: DNS:immich.websec.lab
notBefore=Oct  6 11:27:58 2026   notAfter=Jan  4 11:27:57 2027
# Kette: Leaf <- Pebble Intermediate CA <- Pebble Root CA
\end{lstlisting}

Für die Vertrauensprüfung wird der Webserver (NginX) mit dem ausgestellten
Zertifikat betrieben. Ein Client, der die zugehörige Root-CA kennt, akzeptiert
die Verbindung ohne Warnung; ohne sie schlägt die Prüfung fehl:

\begin{lstlisting}[language=bash]
curl               https://immich.websec.lab/  -> ssl_verify_result=20 (unbekannte CA)
curl --cacert root https://immich.websec.lab/  -> ssl_verify_result=0, HTTP 200
\end{lstlisting}

Das ist exakt der Mechanismus, den auch Let's Encrypt nutzt: Ein über ACME
ausgestelltes, von einer vertrauenswürdigen CA signiertes Zertifikat macht die
Verbindung für den Browser gültig und schützt, wie oben gezeigt, gegen das
Mitlesen durch den MITM-Angreifer. Das Setup liegt reproduzierbar unter
\texttt{pebble-lab/} (\texttt{run-acme.sh}).


\end{document}
